%! Matrix Multiplication and A = CR — Unit 1, Lesson 1.4 — Slides (Beamer)
\documentclass[aspectratio=169, 11pt]{beamer}
\usepackage{linalg-beamer}   % provides \burgundyheader{} and \sectionlabel[color]{}
\newcommand{\xx}{\mathbf{x}}
\newcommand{\bb}{\mathbf{b}}

\begin{document}

% TITLE SLIDE (CourseName is not defined in beamer, so the name is literal)
{
\setbeamercolor{background canvas}{bg=burgundy}
\begin{frame}[plain]
  \vspace{2.8cm}\hspace{0.55cm}%
  \begin{minipage}{0.9\textwidth}
    {\color{goldacc}\bfseries\small LESSON 1.4}\\[0.5cm]
    {\color{white}\bfseries\LARGE Matrix Multiplication and $A=CR$}\\[0.3cm]
    {\color{white}\bfseries\large Keep the columns that are really new --- record a recipe for the rest.}\\[1.0cm]
    {\color{blushmid}\small Linear Algebra~$\cdot$~Unit 1: Vectors and Matrices}
  \end{minipage}
\end{frame}
}

% HOOK
\begin{frame}[t, plain]
  \burgundyheader{Is any column just a mix of the others?}
  \sectionlabel{Hook}
  \begin{columns}[T]
    \begin{column}{0.54\textwidth}
      Read $A=\begin{bsmallmatrix} 1 & 3 & 4 \\ 2 & 1 & 3 \end{bsmallmatrix}$ as three columns
      $\begin{bsmallmatrix} 1\\2 \end{bsmallmatrix},\begin{bsmallmatrix} 3\\1 \end{bsmallmatrix},\begin{bsmallmatrix} 4\\3 \end{bsmallmatrix}$.
      \[
        \begin{bsmallmatrix} 1\\2 \end{bsmallmatrix}+\begin{bsmallmatrix} 3\\1 \end{bsmallmatrix}
        =\begin{bsmallmatrix} 4\\3 \end{bsmallmatrix}.
      \]
      Column~3 is a \textbf{repeat} --- only two columns are truly new.
    \end{column}
    \begin{column}{0.44\textwidth}
      \centering
      \begin{tikzpicture}[scale=0.5]
        \draw[step=1, black!15, thin] (-0.3,-0.3) grid (4.3,3.3);
        \draw[->, thick, black!60] (-0.4,0)--(4.5,0);
        \draw[->, thick, black!60] (0,-0.4)--(0,3.5);
        \draw[->, very thick, burgundy] (0,0)--(1,2) node[above left]{\scriptsize $\mathbf{a}_1$};
        \draw[->, very thick, royalblue] (0,0)--(3,1) node[below right]{\scriptsize $\mathbf{a}_2$};
        \draw[->, very thick, black!70, dashed] (0,0)--(4,3) node[above right]{\scriptsize $\mathbf{a}_3$};
      \end{tikzpicture}
    \end{column}
  \end{columns}
\end{frame}

% MULTIPLY BY COLUMNS
\begin{frame}[t, plain]
  \burgundyheader{Multiply one column at a time}
  \sectionlabel{Definition}
  Each column of $B$ combines the columns of $A$: column $j$ of $AB$ is $A\mathbf{b}_j$.
  \[
    AB = A\begin{bmatrix} \mathbf{b}_1 & \mathbf{b}_2 \end{bmatrix}
    = \begin{bmatrix} A\mathbf{b}_1 & A\mathbf{b}_2 \end{bmatrix}.
  \]
  \begin{itemize}
    \item Example: $\begin{bsmallmatrix} 1 & 3 \\ 2 & 1 \end{bsmallmatrix}\begin{bsmallmatrix} 2 & 1 \\ 1 & 0 \end{bsmallmatrix}
      =\begin{bsmallmatrix} 5 & 1 \\ 5 & 2 \end{bsmallmatrix}$ --- each column is a combination of $\begin{bsmallmatrix} 1\\2 \end{bsmallmatrix},\begin{bsmallmatrix} 3\\1 \end{bsmallmatrix}$.
    \item This is Lesson 1.3, run once per column.
  \end{itemize}
\end{frame}

% A = CR
\begin{frame}[t, plain]
  \burgundyheader{$A=CR$: keep $C$, record the recipe $R$}
  \sectionlabel{Factorization}
  \begin{columns}[T]
    \begin{column}{0.54\textwidth}
      \textbf{$C$} = the independent columns (the truly new ones).\\[4pt]
      \textbf{$R$} = the recipe: each column says how to mix $C$ to rebuild that column of $A$.
      \[
        \underbrace{\begin{bsmallmatrix} 1 & 3 & 4 \\ 2 & 1 & 3 \end{bsmallmatrix}}_{A}
        =\underbrace{\begin{bsmallmatrix} 1 & 3 \\ 2 & 1 \end{bsmallmatrix}}_{C}
         \underbrace{\begin{bsmallmatrix} 1 & 0 & 1 \\ 0 & 1 & 1 \end{bsmallmatrix}}_{R}
      \]
    \end{column}
    \begin{column}{0.42\textwidth}
      \begin{block}{Read $R$}
        Columns 1--2: the identity (the base columns themselves). Column 3: $\begin{bsmallmatrix} 1\\1 \end{bsmallmatrix}$ --- ``one of each,'' rebuilding $\mathbf{a}_3=\mathbf{a}_1+\mathbf{a}_2$.
      \end{block}
    \end{column}
  \end{columns}
\end{frame}

% RANK
\begin{frame}[t, plain]
  \burgundyheader{Rank = how many columns are really new}
  \sectionlabel[goldacc]{Rank}
  \begin{block}{The count}
    The \textbf{rank} is the number of independent columns --- the number of columns in $C$.
    Independent columns fill the \textbf{plane} (rank 2); parallel columns fill only a \textbf{line} (rank 1).
  \end{block}
  For $\begin{bsmallmatrix} 1 & 3 & 4 \\ 2 & 1 & 3 \end{bsmallmatrix}$ the rank is $2$. For
  $\begin{bsmallmatrix} 1 & 2 \\ 2 & 4 \end{bsmallmatrix}$ (parallel columns) the rank is $1$: $C$ is one column, $R$ one row.
\end{frame}

% WRAP / PREVIEW
\begin{frame}[t, plain]
  \burgundyheader{Wrap-up and what's next}
  \sectionlabel{Connections}
  \textbf{Today:} multiply by columns ($AB=\begin{bsmallmatrix} A\mathbf{b}_1 & A\mathbf{b}_2 \end{bsmallmatrix}$);
  factor $A=CR$ with $C$ the independent columns and $R$ the recipe; the rank counts the columns of $C$.\\[8pt]
  \textbf{Next --- Unit 2:} solving $A\xx=\bb$ --- given a target $\bb$, exactly when can we hit it,
  and how do we find the weights?
\end{frame}

\end{document}
